\section{Related Work}
\label{sec:related}

\subsection{Vision-based fall detection}
Fall detection has been studied for two decades, mostly for older people living alone \cite{mubashir2013survey,igual2013challenges}, and camera-based methods have been reviewed recently in \cite{gutierrez2021review,alam2022review}. Early methods used hand-crafted cues such as the aspect ratio of a person's silhouette, the orientation of a fitted ellipse and the height of the body above the floor; UR Fall's own baseline, for example, classifies depth-map features including the ratio of current to standing height \cite{kwolek2014urfall}. Deep networks trained on fall videos later replaced most of these features \cite{alam2022review}. A third line of work runs a pose estimator such as OpenPose \cite{cao2021openpose} or YOLO-Pose \cite{maji2022yolopose} and classifies the keypoints, either with learned models \cite{ramirez2021skeleton} or with rules on the speed of the hip centre, the angle of the body's centre line and the width-to-height ratio \cite{chen2020openposefall}. Our baseline follows the rule-based variant. Detector-based systems aimed at industrial sites also exist \cite{pereira2024industrialfall}.

Most of these methods are trained and tested on a few staged datasets \cite{kwolek2014urfall,martinez2019upfall,auvinet2010mcfd,charfi2013le2i}. OmniFall \cite{schneider2025omnifall} unifies eight of them under one 16-class taxonomy with cross-subject and cross-view splits, and adds synthetic and real-accident test data. Its results are directly relevant here: zero-shot multimodal models detected genuine falls about as well as fine-tuned models, but kept confusing a person who has \textit{fallen} with a person who is \textit{lying} down. Published results rarely break false alarms down by the posture that caused them; we do.

\subsection{Vision--language models and small VLMs}
Contrastive image--text pre-training \cite{radford2021clip} made zero-shot image classification practical, and connecting a vision encoder to a language model produced generative VLMs that answer free-form questions about images \cite{li2023blip2,liu2023llava}. Later work has pushed these models down to a few billion parameters or fewer for phones and embedded boards \cite{chu2023mobilevlm,beyer2024paligemma,marafioti2025smolvlm,bai2025qwen25vl,amini2025lfm2,qwen2026qwen35}, often with vision encoders that accept native aspect ratios and variable resolution \cite{tschannen2025siglip2}. The LFM2 family that we study was designed by searching architectures under on-device latency and memory constraints \cite{amini2025lfm2}. Model releases usually report general benchmarks and throughput on desktop or phone hardware; they say little about behaviour on one safety-critical yes/no question, or about cost on a particular edge board.

VLMs have known weaknesses on yes/no questions. They often answer ``yes'' when asked about objects that are not in the image \cite{li2023pope}, and this yes-bias is common enough to be studied as a failure mode in its own right \cite{guan2024hallusionbench,chan2026yesbias}. Their answers also change with superficial changes to the prompt \cite{zong2024permutations}. The bias does not always point the same way: in zero-shot surveillance anomaly detection, Yao et al.\ found that multimodal models favoured the ``normal'' class so strongly that recall collapsed, and that class-specific instructions moved the decision boundary \cite{yao2026surveillance}. For a \pd{} detector either direction is costly, as missed events or as false alarms, which is why we vary the prompt and report operating points chosen on data.

\subsection{VLMs for fall detection and workplace safety}
Few studies apply VLMs to fall detection. Junca et al.\ used a production multimodal model as the decision stage of a thermal-camera fall screener \cite{junca2026thermal}. They found that a single cautionary sentence in the prompt cost a large share of the detected falls, that identical requests at temperature zero still disagreed on 4\% of clips, and that removing the sentence led the model to call deliberate lie-downs falls, the same ambiguity our labelling policy addresses. Shahabian et al.\ compared zero-shot prompting and a few-shot probe on a VLM with skeleton-based and rule-based detectors for an assistive robot \cite{shahabian2026vlmfall}. OmniFall evaluated large zero-shot models on staged, synthetic and real falls \cite{schneider2025omnifall}.

In industrial safety, VLMs have been used to check personal protective equipment \cite{chen2024clip2safety}, detect safety violations in mining video \cite{wu2025monitorvlm}, recognise construction workers' actions \cite{bui2026construction} and inspect construction sites \cite{chen2026constructionsite10k}. SteelBench evaluates VLMs on real steel-plant CCTV, where workers appear small and are often occluded, and finds them well below human performance \cite{yarrabothula2026steelbench}. Closest to our setting, Adil et al.\ evaluated six VLMs under 4\,B parameters zero-shot for construction hazards and improved them by feeding object-detector output into the prompt \cite{adil2026svlm}.

These studies either use large or cloud-hosted models, or evaluate small models on accuracy alone. We found no study that measures how quantization, input resolution and temporal context trade accuracy against latency, memory and power for \pd{} detection on edge hardware, or that reports false alarms per look-alike posture. That is the gap this paper addresses.

\subsection{Quantization and on-device inference}
Post-training quantization reduces language-model weights to 8 or 4 bits with small accuracy loss \cite{dettmers2022llmint8,frantar2023gptq,lin2024awq}. llama.cpp \cite{gerganov2023llamacpp} implements block-wise integer formats in its GGUF file format and runs on CPUs and on Apple, NVIDIA and other GPUs, which makes it a common runtime for small models on edge devices. Laskaridis et al.\ benchmarked language models on phones and Jetson boards, including energy use and the accuracy cost of quantization \cite{laskaridis2024melt}, and Taherin et al.\ separated prefill from decode costs for llama.cpp and Transformers on several Jetson generations \cite{taherin2026hydra}. MLPerf defines standard latency and throughput scenarios for inference benchmarks \cite{reddi2020mlperf}. We follow this practice of reporting tail latency under a fixed scenario, and add two factors that matter for vision: the cost of encoding the image, which dominates latency when the answer is one word, and the effect of image resolution on both cost and accuracy.
